%!TEX root = ../template.tex
%%%%%%%%%%%%%%%%%%%%%%%%%%%%%%%%%%%%%%%%%%%%%%%%%%%%%%%%%%%%%%%%%%%%
%% appendix2.tex
%% NOVA thesis document file
%%
%% Chapter with example of appendix with a short dummy text
%%%%%%%%%%%%%%%%%%%%%%%%%%%%%%%%%%%%%%%%%%%%%%%%%%%%%%%%%%%%%%%%%%%%

\typeout{NT FILE appendix2.tex}%

\chapter{Prompts used for Expirements}
\label{app:prompt_testing}

\section{Zero-shot Prompt}
\begin{lstlisting}[caption={Zero-shot prompt}, label={lst:zero_shot},basicstyle=\ttfamily\tiny]
You are a software engineer who is creating a < project management system > who has several years of experience in created class diagram models from requirements specification.
Based on UML rules and principles, create a class diagram model for this requirements: 
    The project manager leads a team to execute the project within the project's start and end dates.
    Once the project is created in the project management system, a manager may initiate and later terminate the project due to its completion or for some other reason.
    The team executes the project.
    Requirements is the input of a project.
    As outputs the project shall produce a system.
    Requirements and systems are work products.
    Work product has a description, information of the percentage of work completed, and may be validated.
    The validation is depending on the type of work product.
    The requirements are validated with users in workshops.
    System are validated by being tested against the requirements.
    The requirements may be published using various types of media, including on an intranet or paper form.
    Systems may be deployed onto specific plataforms. 

The output must be presented using PlantUML code to create the class diagram model.


Relations between classes:
Type	Symbol	Purpose
Extension	<|--	Specialization of a class in a hierarchy
Implementation	<|..	Realization of an interface by a class
Composition	*--	The part cannot exist without the whole
Aggregation	o--	The part can exist independently of the whole
Dependency	-->	The object uses another object
Dependency	..>	A weaker form of dependency

Declaring element:
    @startuml
    abstract        abstract
    abstract class  "abstract class"
    annotation      annotation
    circle          circle
    ()              circle_short_form
    class           class
    class           class_stereo  <<stereotype>>
    diamond         diamond
    <>              diamond_short_form
    entity          entity
    enum            enum
    exception       exception
    interface       interface
    metaclass       metaclass
    protocol        protocol
    stereotype      stereotype
    struct          struct
    @enduml

Label on relations:
    @startuml

    Class01 "1" *-- "many" Class02 : contains

    Class03 o-- Class04 : aggregation

    Class05 --> "1" Class06

    @enduml

    @startuml
    class Car

    Driver - Car : drives >
    Car *- Wheel : have 4 >
    Car -- Person : < owns

    @enduml


\end{lstlisting}

\section{Few-shot Prompt}
\begin{lstlisting}[caption={Few-shot prompt}, label={lst:Few_shot},basicstyle=\ttfamily\tiny]
You are a software engineer who is creating a < project management system > who has several years of experience in created class diagram models from requirements specification.
Based on UML rules and principles, create a class diagram model for this requirements: 
    The project manager leads a team to execute the project within the project's start and end dates.
    Once the project is created in the project management system, a manager may initiate and later terminate the project due to its completion or for some other reason.
    The team executes the project.
    Requirements is the input of a project.
    As outputs the project shall produce a system.
    Requirements and systems are work products.
    Work product has a description, information of the percentage of work completed, and may be validated.
    The validation is depending on the type of work product.
    The requirements are validated with users in workshops.
    System are validated by being tested against the requirements.
    The requirements may be published using various types of media, including on an intranet or paper form.
    Systems may be deployed onto specific plataforms. 

The output must be presented using PlantUML code to create the class diagram model.

Here are some examples of how you should create the class diagram class:

Input Requirements: 
    Build a software application to manage filmed scenes for realizing a movie, by following the "Hollywood Approach".
    Every scene shall be identified by a code and it is described by a text in natural language.
    Every scene is filmed at least from one different position.
    Each position is a setup.
    Setup is characterized by a code (a string) and a text in natural language where the photographic parameters are noted (e.g., aperture, exposure, focal length, filters, etc.)
    Each setup is related to a single scene 
    Every setup has at least one take filmed in.
    A take has a positive natural number, a real number representing the number of meters of film that have been used for shooting the take, and the code (string) of the reel where the film is stored.
    A take is associated to a single setup.
    Scenes can be internals, filmed in a theater, or a externals, filmed in a location and can either be "day scene" or "night scene".
    Location have a code (string), an address of the location and a text describing the location.
Output Class diagram PlantUML code:
        @startuml
            class Take{
                nbr: integer
                filmed_meters : Real
                reel : string
            }
            class Setup{
                code: string
                photographic_pars: Text
            }
            class Scene{
                code: string
                description: Text
            }
            class Internal{
                theater: string
            }
            class External{
                night_scene : boolean
            }
            class Location{
                name: string
                address: string
                description: text
            }
            Take "1..*" -- "1" Setup: tk_of_stp
            Setup "1..*" -- "1" Scene: stp_for_scn
            Scene <|-- Internal
            Scene <|-- External
            External "0..*" -- "1" Location : located
        @enduml

Input Requirements:
    A user can open a new or existing document.
    Text is entered through a keyboard.
    A document is made up of several pages.
    Each page is made up of a header, body and footer.
    Date, time and page number may be added to header or footer.
    Document body is made up of sentences, which themselves made up of words, and punctuation characters.
    Words are made up of letters, digits and/or special characters.
    Pictures and tables may be inserted into the document body.
    Tables are made up of rows and columns and every cell in a table can contain both text and pictures.
    Users can save or print documents.
Output Class diagram PlantUML code:
    @startuml
        class Document{
            - numberOfPages
            + open()
            + save()
            + print()
            + new()
        }
        class Page{
            - pageNumber
            + newPage()
            + hideHeader()
            + hideFooter()
            + insertPicture()
            + insertTable()
        }
        class Trimming{
            - date
            - time
            - pageNumber
            + display()
            + change()
            + hide()
        }
        class Character{
            - ASCIICode
            - type
            + normal()
            + italic()
            + bold()
            + underline()
        }
        class Table{
            - rows
            - columns
            + insertRow()
            + insertColumn()
            + newTable()
            + insertPicture()
        }
        class Picture{
            - imageType
            - imageSize
            + insert()
            + delete()
        }
        class Cell{
            - cellFormat
            + edit()
        }
        Document "1" *-- "1..*" Page
        Page "1" o--  "0..1" Trimming
        Page "1" o--  "0..*" Character
        Page "1" o--  "0..*" Table
        Page "1" o--  "0..*" Picture
        Character "0..*" --o "1" Cell
        Table "1" *-- "1..*" Cell  
        Picture "0..*" --o "1" Cell
    @enduml

Relations between classes:
Type	Symbol	Purpose
Extension	<|--	Specialization of a class in a hierarchy
Implementation	<|..	Realization of an interface by a class
Composition	*--	The part cannot exist without the whole
Aggregation	o--	The part can exist independently of the whole
Dependency	-->	The object uses another object
Dependency	..>	A weaker form of dependency

Declaring element:
    @startuml
    abstract        abstract
    abstract class  "abstract class"
    annotation      annotation
    circle          circle
    ()              circle_short_form
    class           class
    class           class_stereo  <<stereotype>>
    diamond         diamond
    <>              diamond_short_form
    entity          entity
    enum            enum
    exception       exception
    interface       interface
    metaclass       metaclass
    protocol        protocol
    stereotype      stereotype
    struct          struct
    @enduml

Label on relations:
    @startuml

    Class01 "1" *-- "many" Class02 : contains

    Class03 o-- Class04 : aggregation

    Class05 --> "1" Class06

    @enduml

    @startuml
    class Car

    Driver - Car : drives >
    Car *- Wheel : have 4 >
    Car -- Person : < owns

    @enduml



\end{lstlisting}

\section{Few-shot Chain-of-thought  Prompt}
\begin{lstlisting}[caption={Few-shot Chain-of-thought  Prompt}, label={lst:Few_shot_Chain_of_thought_Prompt},basicstyle=\ttfamily\tiny]


    You are an expert software architect and UML modeling specialist with extensive experience in transforming system requirements into precise PlantUML class diagrams that are blueprints of the system. Your task is to analyze system descriptions and create comprehensive, syntactically correct class diagrams that accurately represent the system's structure and relationships.

    ## Step by Step to create a class diagram

    Before creating the class diagram, follow this systematic analysis:

    ### Step 1: Entity Identification
    - **Identify Nouns**: Extract all significant nouns from the requirements (these become potential classes)
    - **Classify Entities**: Categorize as concrete classes, abstract classes, or interfaces
    - **Eliminate Redundancies**: Remove duplicate or overly similar concepts

    ### Step 2: Relationship Analysis
    - **Identify Associations**: Look for "has-a", "uses", "manages", "contains" relationships
    - **Determine Inheritance**: Find "is-a" relationships and generalization hierarchies
    - **Establish Multiplicity**: Analyze quantitative relationships.Always specify multiplicities (e.g., `"1" *-- "1..*"`)
    - **Identify Aggregation/Composition**: Determine part-whole relationships and their strength

    ### Step 3: Attribute and Method Extraction
    - **Extract Attributes**: Identify properties, characteristics, and data each class should contain
    - **Determine Methods**: Identify behaviors, operations, and responsibilities
    - **Apply Encapsulation**: Consider visibility modifiers (private, protected, public)

    ### Step 4: UML Principles Application
    - **Single Responsibility**: Ensure each class has one clear purpose
    - **Cohesion**: Group related attributes and methods together
    - **Coupling**: Minimize dependencies between classes
    - **Abstraction**: Use abstract classes/interfaces where appropriate

    ### Step 5: Validate & Refine
    - **No plantUML syntax errors** : Ensure that there are no syntax errors in the generated plantUML code.
    - **All requirements mapped** : Guarantee that all specifications of the system present in the description are represented.

    ## PlantUML Syntax Guidelines

    ### Class Declaration
    ```plantuml
    class ClassName {
        - privateAttribute: Type
        # protectedAttribute: Type
        + publicAttribute: Type
        --
        - privateMethod(): ReturnType
        + publicMethod(param: Type): ReturnType
        {abstract} abstractMethod(): ReturnType
        {static} staticMethod(): Type
    }

    abstract class AbstractClassName
    interface InterfaceName
    ```

    ### Relationship Syntax
    - **Association**: `ClassA -- ClassB : label`
    - **Directed Association**: `ClassA --> ClassB : label`
    - **Aggregation**: `ClassA o-- ClassB : label`
    - **Composition**: `ClassA *-- ClassB : label`
    - **Inheritance**: `ChildClass --|> ParentClass`
    - **Interface Implementation**: `Class ..|> Interface`
    - **Dependency**: `ClassA ..> ClassB : uses`

    ### Multiplicity Notation
    - `"1"` - exactly one
    - `"0..1"` - zero or one
    - `"1..*"` - one or more
    - `"0..*"` - zero or more
    - `"n"` - exactly n
    - `"m..n"` - between m and n

    ### Advanced PlantUML Features
    - **Packages**: `package PackageName { ... }`
    - **Stereotypes**: `<<stereotype>>`


    ## Output Requirements

    1. **Structure**: Begin with `@startuml` and end with `@enduml`
    2. **Naming**: Use PascalCase for classes, camelCase for attributes/methods
    3. **Organization**: Group related classes logically
    4. **Clarity**: Include meaningful relationship labels and multiplicities
    5. **Completeness**: Represent all significant entities and relationships from requirements

    ## Examples

    ### Example 1: Movie Scene Management System

    **Input System Description:**
    We are interested in building a software application to manage filmed scenes for realizing a movie, by following the so-called "Hollywood Approach". Every scene is identified by a code (a string) and it is described by a text in natural language. Every scene is filmed from different positions (at least one), each of this is called a setup. Every setup is characterized by a code (a string) and a text in natural language where the photographic parameters are noted (e.g., aperture, exposure, focal length, filters, etc.). Note that a setup is related to a single scene. For every setup, several takes may be filmed (at least one). Every take is characterized by a (positive) natural number, a real number representing the number of meters of film that have been used for shooting the take, and the code (a string) of the reel where the film is stored. Note that a take is associated to a single setup. Scenes are divided into internals that are filmed in a theater, and externals that are filmed in a location and can either be "day scene" or "night scene". Locations are characterized by a code (a string) and the address of the location, and a text describing them in natural language.

    **Analysis Process:**
    1. **Entities**: Scene, Setup, Take, Internal, External, Location
    2. **Relationships**: Scene-Setup (1:many), Setup-Take (1:many), External-Location (many:1)
    3. **Inheritance**: Internal and External inherit from Scene
    4. **Attributes**: Extracted from descriptive text

    **Output Class Diagram:**
    ```plantuml
    @startuml
        class Scene {
            - code: string
            - description: text
            + getCode(): string
            + getDescription(): text
        }
        
        class Setup {
            - code: string
            - photographicParameters: text
            + getPhotographicParameters(): text
            + addTake(take: Take): void
        }
        
        class Take {
            - number: integer
            - filmedMeters: real
            - reelCode: string
            + getFilmedMeters(): real
            + getReelCode(): string
        }
        
        class Internal {
            - theater: string
            + getTheater(): string
        }
        
        class External {
            - isDayScene: boolean
            + isDayScene(): boolean
        }
        
        class Location {
            - code: string
            - address: string
            - description: text
            + getAddress(): string
        }
        
        Scene <|-- Internal
        Scene <|-- External
        Scene "1" *-- "1..*" Setup : contains
        Setup "1" *-- "1..*" Take : has
        External "0..*" --> "1" Location : filmed_at
    @enduml
    ```


    ## Task Execution

    Now apply this framework to analyze the given system requirements and create a comprehensive PlantUML class diagram. Follow the chain of thought process, explicitly showing your analysis before presenting the final diagram.

    ### System Description
    A project manager uses the project management system to manage a project. The project manager leads a team to execute the project within the project's start and end dates. Once a project is created in the project management system, a manager may initiate and later terminate the project due to its completion or for some other reason. As input, a project uses requirements. As output, a project produces a system (or part of a system). The requirements and system are work products: things that are created, used, updated, and elaborated on throughout a project. Every work product has a description, is of some percent complete throughout the effort, and may be validated. However, validation is dependent on the type of work product. For example, the requirements are validated with users in workshops, and the system is validated by being tested against the requirements. Furthermore, requirements may be published using various types of media, including on an intranet or in paper form; and systems may be deployed onto specific platforms.

\end{lstlisting}
